% !TeX program = xelatex
% ============================================================
% interim_report/interim_report.tex — «Промежуточный отчёт» исследовательского
% курсового проекта ОП «Программная инженерия» — деливерабл КТ1.
%
% Структура задана §2.3.2.4 Программы практики (отчётная документация к КТ1
% исследовательского проекта): титульный лист (с УДК) → содержание → основные
% термины, определения и сокращения → введение (актуальность, объект и предмет,
% методы, цели и задачи; новизна и достоверность ожидаемых результатов, их
% теоретическая значимость и практическая ценность) → обзор и анализ источников,
% выбор методов/алгоритмов/моделей → список использованных источников (ГОСТ) →
% Приложение — календарный план работ с указанием дат.
%
% Это НЕ финальный Отчёт (ГОСТ 7.32-2017): реферата, экспериментальной части
% и заключения здесь нет — они появятся в документе «Отчёт» к защите (КТ3).
% Сдаётся на КТ1 без подписей (зачёт/незачёт). На КТ2 промежуточный отчёт
% повторно не сдаётся (там — ссылка на исходный код + актуализированное
% описание проекта).
%
% Как пользоваться: пишите свой текст вместо жёлтых \hseFill{…}; данные о
% себе, руководителе, группе, годе и УДК задайте в настройках проекта.
% ============================================================
\documentclass[12pt,a4paper]{extreport}
\input{preamble}
\begin{document}
\input{title}

% ── Содержание ──
\tableofcontents
\clearpage

% ============================================================
\chapter*{ОСНОВНЫЕ ТЕРМИНЫ, ОПРЕДЕЛЕНИЯ И СОКРАЩЕНИЯ}
\addcontentsline{toc}{chapter}{Основные термины, определения и сокращения}
% ============================================================
В~настоящем отчёте применяют следующие термины с~соответствующими
определениями и сокращения.

\begin{description}[leftmargin=0pt, labelindent=0pt, font=\bfseries, style=nextline]
  \item[\hseFill{Термин 1}] \hseFill{определение}.
  \item[\hseFill{Термин 2}] \hseFill{определение}.
  \item[\hseFill{СОКР}] \hseFill{расшифровка сокращения}.
\end{description}

\clearpage
% ============================================================
\chapter*{ВВЕДЕНИЕ}
\addcontentsline{toc}{chapter}{Введение}
% ============================================================
% §2.3.2.4: актуальность, объект и предмет, методы, цели и задачи;
% новизна и достоверность ожидаемых результатов, их значимость.
\textbf{Актуальность.} \hseFill{какая проблема решается, кого она затрагивает и
почему существующие подходы недостаточны}~\cite{gost732}.

\textbf{Объект исследования} — \hseFill{что изучается в целом}.

\textbf{Предмет исследования} — \hseFill{конкретная сторона объекта}.

\textbf{Методы исследования:} \hseFill{какими методами планируется решать
задачи (анализ источников, математическое моделирование, вычислительный
эксперимент, эмпирическая оценка и т.\,п.)}.

\textbf{Цель работы} — \hseFill{одна формулировка достигаемого научного результата}.

\textbf{Задачи:}
\begin{enumerate}[label=\arabic*), leftmargin=1.25cm]
  \item \hseFill{провести обзор существующих подходов};
  \item \hseFill{предложить и формализовать метод (модель, алгоритм)};
  \item \hseFill{реализовать программное средство и провести эксперимент};
  \item \hseFill{проанализировать результаты и сделать выводы}.
\end{enumerate}

\textbf{Новизна и достоверность ожидаемых результатов.}
\hseFill{в чём новизна ожидаемых результатов и за счёт чего будет обеспечена
их достоверность (теоретическое обоснование, воспроизводимый эксперимент,
сравнение с базовыми методами)}.

\textbf{Теоретическая значимость и практическая ценность.}
\hseFill{что ожидаемые результаты дадут теории и где могут применяться на
практике}.

\clearpage
% ============================================================
\chapter{ОБЗОР И АНАЛИЗ ИСТОЧНИКОВ, ВЫБОР МЕТОДОВ}
% ============================================================
\section{Обзор и анализ источников}
Дайте аналитический обзор предметной области и существующих подходов со
ссылками на источники~\cite{article, webresource}. На каждый источник из
списка в~тексте обязательно должна быть ссылка.

\section{Сравнительный анализ существующих подходов}
Сопоставьте известные методы по значимым критериям; результат сравнения
приведён в~таблице~\ref{tab:approaches}.

\begin{table}[H]
\caption{Сравнение существующих подходов}
\label{tab:approaches}
\centering
\begin{tabular}{|>{\raggedright\arraybackslash}p{5cm}|c|c|c|}
\hline
\textbf{Критерий} & \textbf{Подход~А} & \textbf{Подход~Б} & \textbf{Предлагаемый} \\
\hline
Критерий~1 & \cmpplus  & \cmpminus  & \cmpplus \\
\hline
Критерий~2 & \cmpminus & \cmpplanned & \cmpplus \\
\hline
\end{tabular}
\end{table}

\section{Выбор методов, алгоритмов и моделей}
Обоснуйте выбор методов, алгоритмов и моделей для решения поставленных
задач: \hseFill{какие методы выбраны и почему именно они подходят для
данной задачи}.

\clearpage
% ── Список источников по ГОСТ Р 7.0.5-2008 (методичка ФКН/ПИ): единый
%    алфавитный нумерованный список, сначала латиница, затем кириллица;
%    разделитель областей «. —»; сетевые ресурсы — «[Электронный ресурс].
%    — URL: … (дата обращения: ДД.ММ.ГГ)». ──
\phantomsection
\addcontentsline{toc}{chapter}{Список использованных источников}
\begin{thebibliography}{99}
  \bibitem{article}
  Фамилия, И.~О. Название статьи / И.~О.~Фамилия // Название журнала. --- 2024. --- № 1. --- С.~10--20.
  \bibitem{webresource}
  Название ресурса [Электронный ресурс]. --- URL: \url{https://example.org} (дата обращения: 01.01.2026).
  \bibitem{gost732}
  ГОСТ 7.32--2017. Отчёт о~научно-исследовательской работе. Структура и~правила оформления. --- М.: Стандартинформ, 2017. --- 27~с.
\end{thebibliography}

\clearpage
% ── Приложение — календарный план работ с указанием дат (§2.3.2.4) ──
\appendix
\renewcommand{\thechapter}{\Asbuk{chapter}}
\titleformat{\chapter}[hang]
  {\normalfont\large\bfseries\centering}{Приложение~\thechapter}{1ex}{}
\chapter{Календарный план работ}
\begin{longtable}{|>{\raggedright\arraybackslash}p{3.2cm}|>{\raggedright\arraybackslash}p{5.2cm}|>{\raggedright\arraybackslash}p{3.6cm}|>{\raggedright\arraybackslash}p{2.6cm}|}
\hline
\textbf{Этап работ} & \textbf{Содержание работ} & \textbf{Ожидаемый результат} & \textbf{Даты} \\ \hline
\endhead
\hseFill{этап 1} & \hseFill{что делается} & \hseFill{результат} & \hseFill{ДД.ММ--ДД.ММ} \\ \hline
\hseFill{этап 2} & \hseFill{что делается} & \hseFill{результат} & \hseFill{ДД.ММ--ДД.ММ} \\ \hline
\hseFill{этап 3} & \hseFill{что делается} & \hseFill{результат} & \hseFill{ДД.ММ--ДД.ММ} \\ \hline
\end{longtable}


\end{document}
